\section{Quadratic twists of the actual cubic strict groups}
\label{sec:mixed-cubic-quadratic-twist-strict}

Put $K=\mathbb Q(\omega)$, $t=\sqrt[3]{269}$, and $L=K(t)$.
Fix either faithful character $\chi$ of
$\operatorname{Gal}(L/K)=C_3$.  For
$d\in\mathcal D:=\{-1,2,-2\}$, write
\[
 \begin{gathered}
 \begin{array}{c|ccc}
 d&-1&2&-2\\\hline
 \delta_d&\chi_{-4}&\chi_8&\chi_{-8}
 \end{array}\\
 L_d=L(\sqrt d),\qquad B_d=\mathbb Q(t,\sqrt d).
 \end{gathered}
 \tag{MCQ.1}
\]
Here the Dirichlet characters are restricted to $G_K$ whenever
used as coefficients.  The exact class inputs proved below are
\[
 \begin{array}{c|ccc}
 d&-1&2&-2\\\hline
 \operatorname{Cl}(B_d)\text{ is a quotient of}&C_7&C_{63}&C_{42}.
 \end{array}
 \tag{MCQ.2}
\]
Use the same explicit coefficient-prime set as in (MCP.1):
\[
 \begin{split}
 \mathscr P=\{\ell\text{ prime}:\ &\ell>3,\quad \ell\ne7,269,\\
 &\ell\equiv1\pmod3,\quad
 269^{(\ell-1)/3}\not\equiv1\pmod\ell\}.
 \end{split}
 \tag{MCQ.3}
\]
The three quotient orders have no prime divisors beyond $2,3,7$,
so they impose no additional exclusions in this set.

For $\ell\in\mathscr P$, let $O$ be the original ring of
integers of a finite extension $E/\mathbb Q_\ell$ containing
$\mu_3$, with uniformizer $\varpi$ and residue field $k$.
Let $\kappa:G_K\to1+\mathfrak m_O$ be continuous and unramified
away $\ell$.  For $j=1,2$ and $e\geq1$ with $\ell-1\mid e$,
put
\[
 D=(E/O)(\chi^j\delta_d\kappa^{-1}c_\ell^{-e}).
 \tag{MCQ.4}
\]
The full strict condition means zero restriction to
$H^1(K_v,D)$ at every $v\mid\ell$, and to $H^1(I_v,D)$
at every other finite place; use the same conditions for
$D[\varpi]$.

\begin{theorem}[AC-G207: original-coefficient strict zero for cubic quadratic twists]
\label{AC-G207}
All the choices above satisfy
\[
       \mathcal S_\ell(K,D[\varpi])=0,
       \qquad \mathcal S_\ell(K,D)=0.          \tag{MCQ.5}
\]
On each actual host (FG.3), for its original coefficient orders
and full local conditions (FC.2), this applies to the six labels
$\chi^j\delta_d$, $j=1,2$, $d\in\mathcal D$, with every
conductor-seven and auxiliary $q,d$ factor trivial.  The second
use of $d$ in the source's auxiliary pair is unrelated to the
fixed squarefree integer in (MCQ.1).

For integer $T\geq2$, let
\[
 \begin{split}
 \mathcal Q(T)=\{(\ell,\chi^j\delta_d,e):\ &d\in\mathcal D,
  \ \ell\in\mathscr P,\ j=1,2,\\
 &1\leq e<T,\ \ell-1\mid e\}.
 \end{split}
 \tag{MCQ.6}
\]
Then its exact number of original triples is
\[
 |\mathcal Q(T)|=
 6\sum_{\substack{\ell\in\mathscr P\\\ell\leq T}}
              \left\lfloor\frac{T-1}{\ell-1}\right\rfloor.
 \tag{MCQ.7}
\]
The primes $13$ and $19$ belong to $\mathscr P$, giving the
explicit disjoint contribution
$6\lfloor(T-1)/12\rfloor+6\lfloor(T-1)/18\rfloor$.
These six finite labels are different from the two untwisted
cubic labels of AC-G203.
For any retained set $J$ of actual triples, the nonnegative
strict-group remainder defined in (CP.3) has the exact identity
\[
             \mathcal R_J=\mathcal R_{J\setminus\mathcal Q(T)}.
 \tag{MCQ.8}
\]
This does not assert that any individual derived index is zero.
\end{theorem}

\begin{proof}
\proofstep{1. The three exact residual splitting fields.}
The splitting field $L$ of $X^3-269$ has Galois group $S_3$
over $\mathbb Q$ and unique quadratic subfield $K$.
Indeed the polynomial is Eisenstein and its discriminant
$-27\cdot269^2$ is not a rational square; its three roots are
$t,t\omega,t\omega^2$.  None of $\mathbb Q(\sqrt d)$ for
$d=-1,2,-2$ equals $K$.  Consequently
\[
 [L_d:\mathbb Q]=12,\qquad
 \operatorname{Gal}(L_d/K)=C_3\times C_2.      \tag{MCQ.9}
\]
The two characters $\chi^j\delta_d$ are faithful on this
product: if a cube root times a sign equals one, each factor
equals one.  Their splitting field over $K$ is exactly $L_d$.

For the smaller computational field, put $u=t+\sqrt d$.
The degree-three and degree-two fields are disjoint, so
$[B_d:\mathbb Q]=6$.  The exact norm polynomial is
\[
 \begin{split}
 Q_d(X)&=\bigl(X^3+3dX-269\bigr)^2-d\bigl(3X^2+d\bigr)^2\\
 &=X^6-3dX^4-538X^3+3d^2X^2-1614dX+72361-d^3.
 \end{split}
 \tag{MCQ.10}
\]
Its six roots are $t\omega^i\pm\sqrt d$, and they are
pairwise distinct.  Equality between roots with different
cubic indices would put $t$ in the degree-four field
$\mathbb Q(\omega,\sqrt d)$, a contradiction; with the same
cubic index their signs must also agree.  Thus $u$ generates
exactly $B_d$, and $Q_d$ is its irreducible degree-six
polynomial.  This identifies the class-certificate field,
not just a signature or a numerical class-group candidate.

The completed root-owned computation applies
$\operatorname{bnfinit}(Q_d,1)$ followed by
$\operatorname{bnfcertify}(B,1)$ to each of these three exact
polynomials.  Its saved complete output is
\[
 \begin{array}{c|c|c|c|c}
 d&B.\mathrm{sign}&B.\mathrm{no}&B.\mathrm{cyc}
                      &\operatorname{bnfcertify}(B,1)\\\hline
 -1&[0,3]&7&[7]&1\\
 2&[2,2]&63&[63]&1\\
 -2&[0,3]&42&[42]&1.
 \end{array}
 \tag{MCQ.10a}
\]
The exact input, complete output, all three saved structures,
and execution receipt are retained.  The successful process
finished at 09:25:57.443623 UTC on 23 November 2026 with
exit code zero and all three completion markers.
The documented flag-one conclusion is precisely that the
actual class group is a quotient of the displayed computed
group.  This proves (MCQ.2), and hence, for every
$\ell\in\mathscr P$,
\[
                \operatorname{Cl}(B_d)\{\ell\}=0.
 \tag{MCQ.11}
\]
Only this smaller-field assertion is used.  No equality of
class numbers, complete unit basis, or vanishing of the whole
$\ell$-class group of $L_d$ is inferred.  The author inspected
this completed root execution; a separate independent peer
replay is recorded under its actual ownership.  The earlier
pending degree-twelve computation is an unused attempt.
The complete PARI manual clause is the explicit software
certificate input, not a new proof of its full algorithm.

\proofstep{2. Residual coefficients and both coefficient places.}
Reduction is injective on $\mu_6$ because $\ell>3$.
The character $\kappa$ reduces to one, and
$\overline c_\ell^{\,e}=1$ because $\ell-1\mid e$.
Therefore
\[
             D[\varpi]=k(\overline\chi^{\,j}\overline\delta_d)
 \tag{MCQ.12}
\]
has exact residual order six and splitting field $L_d$.
The residue field $k$ is that of the original $O$ throughout.

The prime $\ell\equiv1\pmod3$ splits in $K$, and both
completions equal $\mathbb Q_\ell$.  At either place the
polynomial $X^3-269$ has separable irreducible reduction:
the noncube condition in (MCQ.3) excludes a root, and a
reducible cubic must have a root.  Thus the local extension
from $L/K$ is unramified of degree three and
$\overline\chi^{\,j}(\operatorname{Frob}_v)$ has order three.
The quadratic character $\delta_d$ is unramified at $\ell$
and takes value $1$ or $-1$ there.  Hence, for both $v\mid\ell$,
\[
 \overline\chi^{\,j}(\operatorname{Frob}_v)
       \overline\delta_d(\operatorname{Frob}_v)\ne1.
 \tag{MCQ.13}
\]
An element of order three cannot equal a sign in characteristic
$\ell>3$.  This also proves the assertion if either chosen
cubic character or Frobenius convention is inverted.

\proofstep{3. Exact degree-two descent and the required class components.}
Let $a$ generate the cubic subgroup of
$\operatorname{Gal}(L_d/K)$, let $b$ change the sign of
$\sqrt d$, and let $\tau$ fix $t,\sqrt d$ while inverting
$\omega$.  Then $L_d^{\langle\tau\rangle}=B_d$ and
\[
 \tau^2=1,\qquad \tau a\tau^{-1}=a^{-1},\qquad \tau b=b\tau.
 \tag{MCQ.14}
\]
For $d<0$, this $\tau$ need not be complex conjugation in
our chosen embedding.  It is the specified field automorphism.

Put $C=\operatorname{Cl}(L_d)\{\ell\}$.  For the quadratic
extension $L_d/B_d$, extension and norm of ideal classes obey
$N\circ j=2$ and $j\circ N=1+\tau$.
The first identity gives injection on $\ell$-primary classes;
the second, by averaging with $1/2$, identifies
$C^{\tau=1}$ with $\operatorname{Cl}(B_d)\{\ell\}=0$.
Hence $(1+\tau)/2$ has zero image on $C$, and $\tau=-1$
on $C$.  The conjugation relation in (MCQ.14) then implies
that $a=a^{-1}$ on $C$.  Since $a^3=1$, it follows that
$a=1$ on $C$.  In particular both nontrivial cubic components,
including the two with $b$-character $\delta_d$, are zero.
This proves exactly the required component assertion without
assuming $C=0$.

Write $A=D[\varpi]$.  Restriction gives an injection
$H^1(G_K,A)\to H^1(G_{L_d},k)$ because its kernel is
$H^1(C_3\times C_2,A)=0$, by averaging over six.
A restricted continuous cocycle is an additive homomorphism
into $k$.  For a full strict class it kills every finite
inertia group; at places above $\ell$ it kills the full
local group as well.  Since $L_d$ is totally imaginary,
unramified abelian class field theory factors this map
through $\operatorname{Cl}(L_d)$.  Every homomorphism into $k$
factors through its $\ell$-primary part $C$.  It is equivariant
for the original residual character.  In particular,
\[
        f(a x)=\overline\chi^{\,j}(a) f(x).
 \tag{MCQ.15}
\]
But $a$ acts trivially on $C$, whereas
$\overline\chi^{\,j}(a)-1\ne0$ in $k$.
Consequently $f=0$, proving $\mathcal S_\ell(K,A)=0$.
The same conclusion holds if the Artin convention inverts
$a$, because either cubic eigenvalue is nontrivial.

\proofstep{4. Every local invariant quotient on the original module.}
For each $v\mid\ell$, choose a Frobenius lift.  Its scalar
on $D$ has the nontrivial reduction (MCQ.13), including when
$\kappa$ is ramified.  The scalar minus one is an $O$-unit,
which proves $D^{G_{K_v}}=0$ at both coefficient places.

At a finite $v\nmid\ell$, both $c_\ell$ and $\kappa$ are
unramified.  The inertia image of $\chi^j\delta_d$ has order
dividing six, prime to $\ell$.  If it is nontrivial, some
scalar minus one is an $O$-unit, so $D^{I_v}=0$.
If it is trivial, then $D^{I_v}=D$, a $\varpi$-divisible module.
This covers the places over $2,3,269$ with their original
full inertia conditions.  In all cases,
\[
 D^{G_{K_v}}/\varpi D^{G_{K_v}}=0\ (v\mid\ell),
 \qquad D^{I_v}/\varpi D^{I_v}=0\ (v\nmid\ell).
 \tag{MCQ.16}
\]

\proofstep{5. Original-uniformizer lifting and arbitrary torsion.}
Apply cohomology to
$0\to D[\varpi]\to D\xrightarrow{\varpi}D\to0$.
For every global or prescribed local group $G$, the relevant
exact part is
\[
 D^G/\varpi D^G\longrightarrow H^1(G,D[\varpi])
             \longrightarrow H^1(G,D)[\varpi]\longrightarrow0.
 \tag{MCQ.17}
\]
Given a full strict class killed by $\varpi$, choose a global
residual lift.  Its restrictions to the prescribed local
groups map to zero in $H^1(G,D)$.  By (MCQ.16)--(MCQ.17),
those residual restrictions are zero already.  The lift is
therefore full strict and vanishes by Step 3.  Thus the
original strict group has no $\varpi$-torsion.
Every continuous cocycle from the compact profinite group
$G_K$ to the discrete module $D$ has finite image and is
killed by some power of $\varpi$.  Any nonzero strict class
would consequently yield a nonzero class killed by $\varpi$,
a contradiction.  This proves both assertions of (MCQ.5).

\proofstep{6. Actual labels, distinctness, and the exact remainder.}
For the actual source range $h\geq1$, the field (FG.3)
contains $\mathbb Q(\zeta_8)$, hence $\sqrt{-1}$ and
$\sqrt2$.  The three quadratic characters in (MCQ.1) are
therefore actual characters of its two-power factor.
They remain distinct on $G_K$: the biquadratic field
$\mathbb Q(i,\sqrt2)$ intersects $K$ trivially, since its
three quadratic subfields have discriminants $-4,8,-8$.
The cubic factor is independent by the product decomposition
(FC.1).  Consequently all six $\chi^j\delta_d$ are distinct
and are different from the two untwisted cubic labels in
(MCP.4).

Their values lie in the original completion $O_\ell$ of
$\mathbb Z[\zeta_{3\cdot2^h}]$.  Apply (MCQ.5) with
$\kappa=1$ and the canonical identity
$O_\ell\otimes\mathbb Q_\ell/\mathbb Z_\ell=E_\ell/O_\ell$.
Inflation from the actual $G_{K,\Sigma_\ell}$ is injective;
its strict classes retain all prescribed local conditions,
and are inertia-zero outside $\Sigma_\ell$.  They thus
inject into the zero group just proved.  The ring, uniformizer,
and normalizing rank $b_\ell$ remain the original ones.

For each $(d,\ell)$ the number of positive eligible weights
is $\lfloor(T-1)/(\ell-1)\rfloor$, which is zero unless
$\ell\leq T$.  The six labels and the prime coordinates are
distinct, proving (MCQ.7).  Their cubic factor is nontrivial,
so the summands in (CP.3) are the full original $K$-strict
groups, not rational parity components.  Deleting these
zero summands from any $J$ proves (MCQ.8), with no added
coefficient error or independently chosen window.
\end{proof}

\paragraph{Finite-cover dictionary and scope.}
The finite modules $D[\varpi^n]$ are the original finite
locally constant etale sheaves on
$\operatorname{Spec}\mathcal O_K[1/(2\cdot3\cdot269\cdot\ell)]$.
Their residual scalar is trivialized by the displayed
$C_3\times C_2$ cover $L_d/K$.  The strict kernels use the
same completion and inertia restrictions on these original
coefficients.  This is the retained SGA~1, Expose V,
sections 4 and 7 finite-cover/Artin--Tate dictionary; no
fresh archival reading or motivic realization is claimed.
The root proposed these covers and owns their class computations.
The author supplied the degree-six fixed-transposition reduction
so that full degree-twelve class vanishing is not required.
The author supplies the explicit quadratic-twist adaptation
of the independently checked AC-G203 residual-lifting argument.
The complete unramified class-field statement and the
complete Sharifi 1.3.1--1.3.5 proof passage supply the stated
Artin and norm/extension interfaces; the local lifting proof
is written above.  No GRH-dependent class output is used.
The remaining all-prime budget, $\mathcal R_{\rm cub}$,
and H04.4 are separate gates.
